% !PRINTED 360
\weberpage{360}

Dies giebt die Entwickelung des Bruches $m:m_1$ oder $m:n$ in einen \gesperrt{Kettenbruch}. Wir erhalten successive
\weq{5}{
\frac{m}{n} = a_0 + \cfrac{1}{\dfrac{m_1}{m_2}}, \quad
\frac{m}{n} = a_0 + \cfrac{1}{a_1 + \cfrac{1}{\dfrac{m_2}{m_3}}} \cdots,
}
und endlich
\weq{6}{
\frac{m}{n} = a_0 + \cfrac{1}{a_1 + \cfrac{1}{a_2 + \cdots + \cfrac{1}{a_{\nu-1} + \cfrac{1}{m_\nu}}}},
}
wofür wir, des bequemeren Druckes wegen, auch
\weq{7}{
\begin{aligned}
\frac{m}{n} &= \left(a_0, \frac{m_1}{m_2}\right), \quad \frac{m}{n} = \left(a_0, a_1, \frac{m_2}{m_3}\right)\\
\frac{m}{n} &= (a_0, a_1, a_2 \ldots a_{\nu-1}, m_\nu)
\end{aligned}
}
setzen.

Hierzu ist noch Folgendes zu bemerken: Jeder positive rationale Bruch lässt sich in die Form $m:n$ setzen, so dass $m$ und $n$ relativ prim sind, und daher können wir jeden solchen Bruch in einen Kettenbruch von der Form (6) oder (7) entwickeln. Auch für negative rationale Brüche gilt dies, wenn wir für die Zahl $a_0$ auch negative Werthe zulassen. Die $a_1, a_2 \ldots a_{\nu-1}, m_\nu$ sind aber auch dann positive ganze Zahlen. Diese Zahlen sind vollkommen und unzweideutig bestimmt. Nur in einer Beziehung steht uns noch eine Willkür offen.

Es ist nämlich nach der bis jetzt getroffenen Bestimmung der letzte der Theilnenner $m_\nu$ grösser als 1, da wir angenommen haben, dass $m_{\nu+1}$ der erste unter den Resten sein sollte, der gleich 1 ist. Demnach können wir, wenn $a_\nu$ eine positive ganze Zahl ist, entweder
\[
m_\nu = a_\nu,
\]
oder
\[
m_\nu = a_\nu + \frac{1}{1}
\]
setzen und erhalten daher die zwei Darstellungen durch Kettenbrüche:
%CONT>>
